\subsubsection{Testcase Module 2: Quản lý Lộ trình cá nhân}

\hypertarget{target:tc_m2}{}
\begin{landscape}
\small 
\begin{longtable}{|p{2.5cm}|p{4.0cm}|p{6.5cm}|p{4.5cm}|p{4.5cm}|}
    
    % --- PHẦN 1: TIÊU ĐỀ TRANG ĐẦU ---
\caption{Danh sách Test Case chi tiết - Module 2: Quản lý Lộ trình cá nhân} \label{tab:tc_module2} \\
\hline
\textbf{ID} & \textbf{Tên Test Case} & \textbf{Các bước thực hiện} & \textbf{Dữ liệu kiểm thử} & \textbf{Kết quả mong đợi} \\
\hline
\endfirsthead

% --- PHẦN 2: TIÊU ĐỀ CÁC TRANG SAU ---
\multicolumn{5}{c}{{\bfseries \tablename\ \thetable{}: Danh sách Test Case - Module 2 Quản lý Lộ trình cá nhân}} \\
\hline
\textbf{ID} & \textbf{Tên Test Case} & \textbf{Các bước thực hiện} & \textbf{Dữ liệu kiểm thử} & \textbf{Kết quả mong đợi} \\
\hline
\endhead

% --- PHẦN 3: KẾT THÚC BẢNG ---
\hline
\endlastfoot

% ================= NỘI DUNG TEST CASE MODULE 2 =================

% --- LUỒNG CHÍNH (MAIN FLOW) ---
TC\_M2\_01 & Xem danh sách lộ trình (My Trips) & 
1. Đăng nhập hệ thống. \newline 
2. Truy cập menu ``Chuyến đi của tôi''. & 
- Tài khoản: User đã có 5 lộ trình đã lưu. & 
Hiển thị danh sách các lộ trình dưới dạng thẻ (Card) bao gồm: Ảnh bìa, Tên chuyến, Thời gian, Trạng thái (Sắp đi/Đã qua). \\
\hline

TC\_M2\_02 & Xem chi tiết một lộ trình & 
1. Tại màn hình danh sách. \newline 
2. Nhấn vào thẻ lộ trình ``Đà Lạt 2025''. & 
- Lộ trình: Đà Lạt (Public) & 
Chuyển hướng đến trang chi tiết lộ trình với đầy đủ timeline, bản đồ và tổng chi phí dự kiến. \\
\hline

TC\_M2\_03 & Tìm kiếm lộ trình thành công & 
1. Nhập từ khóa vào ô tìm kiếm. \newline 
2. Nhấn Enter. & 
- Keyword: ``Hà Nội'' \newline
- Data: Có 2 lộ trình khớp tên. & 
Danh sách lọc lại, chỉ hiển thị 2 lộ trình có tên chứa từ khóa ``Hà Nội''. \\
\hline

% --- LUỒNG THAY THẾ (ALTERNATIVE FLOW) ---
TC\_M2\_04 & Lọc lộ trình theo trạng thái & 
1. Nhấn vào bộ lọc (Filter). \newline 
2. Chọn trạng thái ``Đã kết thúc''. & 
- Trạng thái: Past Trips & 
Chỉ hiển thị các chuyến đi có ngày kết thúc nhỏ hơn ngày hiện tại. \\
\hline

TC\_M2\_05 & Xóa nhiều lộ trình cùng lúc (Bulk Delete) & 
1. Nhấn nút ``Chọn nhiều''. \newline 
2. Tích chọn 3 lộ trình. \newline 
3. Nhấn icon Thùng rác $\to$ Xác nhận. & 
- Số lượng: 3 items & 
3 lộ trình đã chọn biến mất khỏi danh sách. Hệ thống thông báo ``Đã xóa 3 mục thành công''. \\
\hline

TC\_M2\_06 & Chia sẻ lộ trình cho bạn bè & 
1. Tại thẻ lộ trình, nhấn icon ``Chia sẻ''. \newline 
2. Nhập email người nhận. \newline 
3. Nhấn ``Gửi lời mời''. & 
- Email: friend@gmail.com \newline
- Quyền: Xem (View only) & 
Hệ thống gửi email mời tham gia/xem lộ trình đến địa chỉ đích. Thông báo ``Đã gửi lời mời''. \\
\hline

TC\_M2\_07 & Nhân bản lộ trình (Duplicate) & 
1. Nhấn menu ``Khác'' (3 chấm) trên thẻ lộ trình. \newline 
2. Chọn ``Nhân bản''. & 
- Lộ trình gốc: ``Quy Nhơn'' & 
Tạo ra một bản sao mới có tên ``Copy of Quy Nhơn'' nằm trên cùng danh sách. \\
\hline

% --- LUỒNG NGOẠI LỆ (EXCEPTION FLOW) ---
TC\_M2\_08 & Tìm kiếm không có kết quả & 
1. Nhập từ khóa không tồn tại. \newline 
2. Nhấn Enter. & 
- Keyword: ``@\#\$\%'' & 
Hiển thị thông báo/hình ảnh minh họa: ``Không tìm thấy lộ trình nào phù hợp với từ khóa của bạn''. \\
\hline

TC\_M2\_09 & Xóa lộ trình khi mất kết nối mạng & 
1. Chọn xóa lộ trình. \newline 
2. Ngắt kết nối Internet ngay lúc nhấn xác nhận. & 
- Mạng: Disconnected & 
Hệ thống hiển thị lỗi: ``Kết nối không ổn định, vui lòng thử lại sau''. Lộ trình chưa bị xóa. \\
\hline

TC\_M2\_10 & Danh sách lộ trình trống (User mới) & 
1. Đăng nhập bằng tài khoản mới tạo. \newline 
2. Vào ``Chuyến đi của tôi''. & 
- User: New account & 
Hiển thị màn hình chào mừng (Empty State) với nút kêu gọi hành động: ``Bạn chưa có chuyến đi nào? Tạo ngay!''. \\
\hline

\end{longtable}

\end{landscape}